%! Solving Polynomial Equations by Factoring — Unit 1, Lesson 1.6 — Homework (Answer Key)
\documentclass[10pt]{article}
\usepackage{atda-article}
\usepackage{atda-key}   % pulls in -boxes; do NOT also load -boxes
\newcommand{\TallMath}[1]{$\displaystyle #1\rule[-1.4em]{0pt}{3.2em}$}
\setlength{\workrowsep}{8pt}

\begin{document}

\pageheader{Unit 1, Lesson 1.6}{Homework --- Answer Key}
% NAMESTRIP: no \namedateperiod here — mirrors the blank. Teacher-only prose goes
% in the lesson plan as \begin{teachernote}[Homework], never in a key.

\vspace{0.15in}

\boxguard[16]
\begin{notesbox}{Practice --- Solve Every Equation Completely}
\small
Show your work. \textbf{Standard form first}, then factor completely, then set each factor equal to
zero. List \emph{every} solution --- and never divide both sides by a variable.

\begin{enumerate}[leftmargin=1.6em, itemsep=6pt, topsep=4pt]

\item Solve: $(x+8)(x-5)=0$
\begin{work}
  x+8 &= 0 \quad\Rightarrow\quad x = -8 \\
  x-5 &= 0 \quad\Rightarrow\quad x = 5
\end{work}

\item Solve: $x^2-5x-14=0$
\begin{work}
  (x-7)(x+2) &= 0 \\
  x &= 7 \quad \text{or} \quad x = -2
\end{work}

\item Solve: $x^2+2x-15=0$
\begin{work}
  (x+5)(x-3) &= 0 \\
  x &= -5 \quad \text{or} \quad x = 3
\end{work}

\item Solve: $x^2-121=0$
\begin{work}
  (x+11)(x-11) &= 0 \\
  x &= -11 \quad \text{or} \quad x = 11
\end{work}

\item Solve: $2x^2=10x$
\begin{work}
  2x^2-10x &= 0 \\
  2x(x-5) &= 0 \\
  x &= 0 \quad \text{or} \quad x = 5
\end{work}

\item Solve: $x^2+8x+16=0$ \quad (what happens to the two solutions here?)
\begin{work}
  (x+4)^2 &= 0 \\
  x &= -4 \quad \text{(a double root --- counted twice)}
\end{work}

\item Solve: $x^2+3x=18$
\begin{work}
  x^2+3x-18 &= 0 \\
  (x+6)(x-3) &= 0 \\
  x &= -6 \quad \text{or} \quad x = 3
\end{work}

\item Solve: $x^3-25x=0$
\begin{work}
  x\left(x^2-25\right) &= 0 \\
  x(x+5)(x-5) &= 0 \\
  x &= 0, \quad x = -5, \quad x = 5
\end{work}

\item Solve: $x^3+3x^2-10x=0$
\begin{work}
  x\left(x^2+3x-10\right) &= 0 \\
  x(x+5)(x-2) &= 0 \\
  x &= 0, \quad x = -5, \quad x = 2
\end{work}

\end{enumerate}
\end{notesbox}

\vspace{0.10in}

\boxguard[16]
\begin{scenariobox}[In Context --- The Water Balloon Launch]{royal}
\small
At the end-of-year field day a water balloon is fired from a $96$-foot press box. Its height above
the field $t$ seconds after launch is
\[ h(t) = -16t^2+80t+96 \quad \text{feet.} \]

\begin{enumerate}[leftmargin=1.6em, itemsep=6pt, topsep=4pt, start=10]

\item Factor $h(t)$ completely. Pull the GCF first --- and $-16$ is part of it.
\begin{work}
  h(t) &= -16t^2+80t+96 \\
       &= -16\left(t^2-5t-6\right) \\
       &= -16(t-6)(t+1)
\end{work}

\item Solve $h(t)=0$. Give both solutions, say which one the situation rejects, and state when the
balloon lands.
\begin{work}
  -16(t-6)(t+1) &= 0 \\
  t &= 6 \quad \text{or} \quad t = -1 \\
  t &= -1 \; \text{rejected (time cannot be negative); the balloon lands at } t=6 \text{ s}
\end{work}

\item The balloon is level with the press box, at $96$ feet, twice. Solve $h(t)=96$, then say in
one sentence what each solution means.
\begin{work}
  -16t^2+80t+96 &= 96 \\
  -16t^2+80t &= 0 \\
  -16t(t-5) &= 0 \\
  t &= 0 \quad \text{or} \quad t = 5
\end{work}
\ansline{$t=0$ is the launch from the press box; $t=5$ s is the balloon falling back past it.}

\end{enumerate}
\end{scenariobox}

\vspace{0.10in}

\boxguard[16]
\begin{scenariobox}[A First Look Ahead]{goldacc}
\small
\begin{enumerate}[leftmargin=1.6em, itemsep=6pt, topsep=4pt, start=13]

\item Next lesson works with expressions like \ \TallMath{\frac{x+5}{x^2-4}}.
\begin{enumerate}[label=(\alph*), leftmargin=1.6em, itemsep=4pt, topsep=3pt]
  \item Solve $x^2-4=0$.
\begin{work}
  (x+2)(x-2) &= 0 \\
  x &= -2 \quad \text{or} \quad x = 2
\end{work}
  \item Division by zero is undefined, so those two values are \textbf{excluded} from the
        expression. In one sentence, say why today's method is exactly how you find the excluded
        values of any such expression.

\ansline{Factoring the denominator and setting each factor to zero is what finds every value that makes it zero.}
\end{enumerate}

\end{enumerate}
\end{scenariobox}

\vspace{0.10in}

\boxguard[18]
\begin{extensionbox}
\small
\textbf{Degree four --- count first, then solve.} Lesson 1.5 factored $x^4-16$ completely. Use that
factorization to solve $x^4-16=0$: how many solutions should there be, how many are \textbf{real},
and how many are \textbf{imaginary}? Say how you know.
\begin{work}
  x^4-16 &= 0 \\
  \left(x^2+4\right)\left(x^2-4\right) &= 0 \\
  \left(x^2+4\right)(x+2)(x-2) &= 0 \\
  x &= -2 \quad \text{or} \quad x = 2 \quad \text{(real)} \\
  x^2+4 &= 0 \; \text{has no real solution, so its two solutions are imaginary}
\end{work}
\ansline{Degree $4$, so four solutions: two real ($\pm 2$) and two imaginary, from the sum of squares.}
\end{extensionbox}

\vspace{0.10in}

\begin{remindbox}
\small
\textbf{Next: Lesson 1.7 --- Simplifying Rational Expressions} (\texttt{A2.EO.1a--b}). Item 13 is
the whole bridge: a rational expression is a fraction of polynomials, and the values that break it
are exactly the solutions of ``denominator $=0$'' --- found by factoring, then the zero product
property. Everything you factored in 1.3--1.5 is about to cancel.
\end{remindbox}

\end{document}
